\begin{frame}{그림으로: 확인 문제의 두 상황}
  \begin{center}
  \includegraphics[height=5.0cm]{ch05_1_handcalc_points.png}
  \end{center}
  \vskip0.2em
  \footnotesize
  \begin{itemize}
    \item \textbf{왼쪽}: 원래 6개 점 --- 서포트 벡터(오렌지)는 $(3,3)$,
      $(-3,-3)$뿐.
    \item \textbf{오른쪽}: $(4,3)\to(0.5,0.5)$로 옮기면 마진 안을 침범 ---
      서포트 벡터가 $(0.5,0.5)$, $(-3,-3)$으로 바뀌며 경계 자체가 이동한다.
    \item 마진(연한 띠)의 폭도 함께 좁아진다.
  \end{itemize}
\end{frame}
